\begin{table}[htbp]\centering\small
\caption{Robustness of the ln GACI coefficient: ln CO2/GDP}\label{tab:A1b}
\begin{tabular}{lccccccccccccc}
\toprule
 & base & +K/L & +sea MA & +LSCI 06- & GACI-region$\times$yr & continent$\times$yr & subregion$\times$yr & region trends & tercile$\times$yr & 1996-2009 & 2010-2019 & 1996-2023 (all) & 2010-2023 \\
\midrule
ln GACI & -0.201 & -0.227 & -0.176 & -0.247** & 0.071 & 0.022 & 0.095 & 0.029 & 0.052 & -0.109 & -0.197* & -0.171 & -0.015 \\
 & (0.162) & (0.179) & (0.159) & (0.111) & (0.151) & (0.172) & (0.162) & (0.144) & (0.136) & (0.128) & (0.119) & (0.134) & (0.079) \\
\midrule
Observations & 3,818 & 3,551 & 3,818 & 1,788 & 3,753 & 3,818 & 3,813 & 3,753 & 3,520 & 2,170 & 1,647 & 4,584 & 2,381 \\
\bottomrule
\end{tabular}
\begin{minipage}{0.95\textwidth}\footnotesize Unified sample: 173 countries, 1996--2019, country-years with all Table 1 outcomes and the instrument observed. All denominators from one source (WDI GDP in constant 2015 US\$ = GDP per capita $\times$ population), so log identities hold exactly. Country and year fixed effects; controls ln population, ln GDP per capita and its square. Standard errors clustered by country. *, **, *** : 10, 5, 1\%. K/L from PWT 11. GACI-region = the 7 macro-regions of the GACI panel (AF, AS, EU, LA, ME, NA, SW); continent = 5 UN continents; subregion = 22 UN subregions; region trends = region-specific linear trends. CO2 columns 12--13 extend the sample to 2023 (CEDS ends 2019).\end{minipage}
\end{table}